\chapter{Complexity — Model Invariance}
%%% generated by the blueprint pipeline from TCSlib.Complexity.ClassP.ModelInvariance
%%%%%%%%%%%%%%%%%%%%%%%%%%%%%%%%%%%%%%%%%%%%%%%%%%%%%%%%%%%%%%%%%%%%%%%%%%%%%%%
\section{Overview}

This module records the model-invariance payoff of the Turing-machine robustness
theorems, in the spirit of Arora--Barak \S1.3.1 and \S1.6.1 (``why the model may
not matter''): the class $\mathrm{DTIME}$ of binary languages is invariant under
enlarging the tape alphabet, the alphabet-dependent constant being absorbed into
$\mathrm{DTIME}$'s built-in constant factor, and the class $\mathsf{P}$ is
invariant both under enlarging the alphabet and under restricting to machines
with a single work tape, since the quadratic slowdown of tape reduction
preserves polynomial time bounds.

%%%%%%%%%%%%%%%%%%%%%%%%%%%%%%%%%%%%%%%%%%%%%%%%%%%%%%%%%%%%%%%%%%%%%%%%%%%%%%%
\section{Declarations}

\begin{definition}[Deciding a binary language via a symbol embedding]
\lean{Turing.FinTM.DecidesInTimeVia}
\label{Turing.FinTM.DecidesInTimeVia}
\leanok
\uses{Turing.FinTM, Turing.FinTM.ComputesInTime, Turing.MultiTapeTM.indicator}
Let $\Gamma$ be an alphabet, let $M$ be a deterministic finite Turing machine
over $\Gamma$, let $e : \{0,1\} \hookrightarrow \Gamma$ be an injective
embedding of the binary symbols into $\Gamma$, let $L \subseteq \{0,1\}^{*}$ be
a language over the binary alphabet, and let $T : \mathbb{N} \to \mathbb{N}$ be
a time bound. We say that $M$ decides $L$ via $e$ within time $T$ if for every
binary string $x$, when $M$ is run on the input obtained by applying $e$ to
each symbol of $x$, it halts within $T(|x|)$ steps with output the one-symbol
string $e(b)$, where $b \in \{0,1\}$ is the indicator of whether $x \in L$.
\end{definition}

\begin{theorem}[Alphabet invariance of DTIME]
\lean{Complexity.mem_DTIME_of_decidesInTimeVia}
\label{Complexity.mem_DTIME_of_decidesInTimeVia}
\leanok
\uses{Complexity.DTIME, Turing.FinTM, Turing.FinTM.DecidesInTimeVia, Turing.FinTM.alphabet_reduction, Turing.MultiTapeTM.indicator}
\difficulty{3}
Let $\Gamma$ be a finite alphabet, let
$e : \{0,1\} \hookrightarrow \Gamma$ be an injective embedding of the binary
symbols, let $L \subseteq \{0,1\}^{*}$ be a language over the binary alphabet,
and let $T : \mathbb{N} \to \mathbb{N}$ be a time bound. If some deterministic
finite Turing machine over $\Gamma$ decides $L$ via $e$ within time $T$, then
$L \in \mathrm{DTIME}(n \mapsto T(n) + 1)$; that is, some deterministic finite
Turing machine over the binary alphabet decides $L$ within
$a \cdot (T(n) + 1)$ steps on every input of length $n$, for some constant $a$.
\end{theorem}

\begin{theorem}[Alphabet invariance of P]
\lean{Complexity.mem_P_of_decidesInTimeVia_poly}
\label{Complexity.mem_P_of_decidesInTimeVia_poly}
\leanok
\uses{Complexity.P, Complexity.mem_DTIME_of_decidesInTimeVia, Complexity.mem_P_of_dtime_le, Complexity.succ_pow_le, Turing.FinTM, Turing.FinTM.DecidesInTimeVia}
\difficulty{4}
Let $\Gamma$ be a finite alphabet, let
$e : \{0,1\} \hookrightarrow \Gamma$ be an injective embedding of the binary
symbols, let $L \subseteq \{0,1\}^{*}$ be a language over the binary alphabet,
and let $C$ and $d$ be natural numbers. If some deterministic finite Turing
machine over $\Gamma$ decides $L$ via $e$ within $C \cdot (n+1)^d$ steps on
every input of length $n$, then $L \in \mathsf{P}$. In particular, polynomial-time
decidability over any finite alphabet already places a language in the
binary-alphabet class $\mathsf{P}$.
\end{theorem}

\begin{theorem}[P is exactly one-work-tape polynomial time]
\lean{Complexity.mem_P_iff_one_work_tape}
\label{Complexity.mem_P_iff_one_work_tape}
\leanok
\uses{Complexity.P, Complexity.mem_P_iff, Turing.FinTM, Turing.FinTM.ComputesInTime.mono, Turing.FinTM.DecidesInTime, Turing.FinTM.one_work_tape_binary, Turing.MultiTapeTM.indicator}
\difficulty{4}
A language $L$ over the binary alphabet belongs to $\mathsf{P}$ if and only if
there exist a deterministic finite Turing machine $M$ over $\{0,1\}$ and
natural numbers $C$ and $d$ such that $M$ has exactly one work tape and $M$
decides $L$ within $C \cdot (n+1)^d$ steps on every input of length $n$. Thus
restricting to a single work tape does not change the class $\mathsf{P}$: the
quadratic overhead of tape reduction preserves polynomial time bounds.
\end{theorem}
